% Appendix A of Singular Learning Theory for the Rest of Us: the shared toolkit.
\appendixchapter{A Toolkit}{Integrals, singularities, and samplers}{ch:toolkit}
\begin{refsection}[chapters/17-toolkit.bib]

This appendix collects the tools that several chapters share: Gaussian integrals and Laplace's method; the Gamma function, Mellin transforms and regular variation; analytic functions and the resolution of their singularities; the random process on which Watanabe's theorems rest; divergences and modes of convergence; and the samplers, estimators and intervals behind the book's figures. Short proofs are given in full. Long ones are replaced by a precise statement and a source, and the statement says what the book takes on trust; the two results of this kind that carry the most weight are Hironaka's theorem (\cref{thm:toolkit:hironaka}) and Watanabe's standard form of the log-likelihood ratio (\cref{thm:toolkit:empirical}). Each item ends by saying which chapters use it. The appendix is a reference, not a course: each section can be read when a chapter points to it.

\section{Gaussian integrals and Laplace's method}\label{sec:toolkit:gaussian}

When $f$ has a nondegenerate minimum, the integral $\int g\,e^{-nf}$ is a Gaussian integral to leading order; this is \term{Laplace's method}\index{Laplace's method}. The section ends with Stein's identity and Gaussian smoothing.

\begin{theorem}[Laplace's method]\label{thm:toolkit:laplace}
Let $W\subset\R^d$ be a Borel set, let $f,g\colon W\to\R$ be measurable, and let $w_0$ be an interior point of $W$. Suppose that
\begin{enumerate}
\item $f$ has four and $g$ two continuous derivatives on an open set containing a closed ball $\bar B(w_0,r_0)\subset W$;
\item $\nabla f(w_0)=0$, and the Hessian $H=\nabla^2f(w_0)$ is positive definite;
\item the minimum is separated: $\eta(\delta)=\inf\{f(w)-f(w_0):w\in W,\ \abs{w-w_0}\ge\delta\}>0$ for every $\delta>0$;
\item $\int_W\abs ge^{-n_0f}<\infty$ for some $n_0$.
\end{enumerate}
Then, as $n\to\infty$ through real values,
\begin{equation}\label{eq:toolkit:laplace}
e^{nf(w_0)}\int_Wg(w)e^{-nf(w)}\dd w=\Bigl(\frac{2\pi}n\Bigr)^{d/2}(\det H)^{-1/2}\bigl(g(w_0)+O(n^{-1})\bigr).
\end{equation}
In particular, if $g=\varphi$ with $\varphi(w_0)>0$, then
\[
-\log\int_W\varphi e^{-nf}=nf(w_0)+\frac d2\log\frac n{2\pi}+\frac12\log\det H-\log\varphi(w_0)+O(n^{-1}).
\]
\end{theorem}

Hypothesis (iii) fails when $f$ comes back down to $f(w_0)$ at the boundary of $W$ or at infinity. The proof uses the following facts.

\begin{lemma}[Gaussian integrals, moments, and tails]\label{lem:toolkit:gaussian}
Let $Z\sim\Normal(0,1)$ and $Y\sim\Normal(0,I_d)$.
\begin{enumerate}
\item For a symmetric positive definite $d\times d$ matrix $A$ and $b\in\R^d$,
\[
\int_{\R^d}e^{-\frac12x^\top Ax+b^\top x}\dd x=(2\pi)^{d/2}(\det A)^{-1/2}e^{\frac12b^\top A^{-1}b}.
\]
\item For real $s>-1$, $\E\abs Z^s=2^{s/2}\Gamma(\frac{s+1}2)/\sqrt\pi$, so that $\E Z^{2k}=(2k-1)(2k-3)\cdots1$. Every moment of $\abs Y$ is finite, and every odd polynomial in $Y$ has mean zero.
\item For $t>0$, $\frac t{1+t^2}\frac{e^{-t^2/2}}{\sqrt{2\pi}}\le\Prob(Z>t)\le\min\bigl(\frac{e^{-t^2/2}}{t\sqrt{2\pi}},\frac12e^{-t^2/2}\bigr)$.
\item For $r\ge0$, $\Prob(\abs Y\ge r)\le2^{d/2}e^{-r^2/4}$.
\end{enumerate}
\end{lemma}

\begin{proof}
The plan is to reduce everything to the one-dimensional integral of $e^{-t^2/2}$. (i) By Fubini's theorem and polar coordinates, the square of $\int_\R e^{-t^2/2}\dd t$ is $\int_0^{2\pi}\!\int_0^\infty e^{-r^2/2}r\dd r\dd\theta=2\pi$. In general, completing the square gives $-\frac12x^\top Ax+b^\top x=-\frac12(x-c)^\top A(x-c)+\frac12b^\top A^{-1}b$ with $c=A^{-1}b$; translating by $c$ and substituting $x=A^{-1/2}y$, whose Jacobian is $(\det A)^{-1/2}$, leaves $\int e^{-\abs y^2/2}\dd y=(2\pi)^{d/2}$.
(ii) By symmetry and the substitution $u=t^2/2$, $\E\abs Z^s=\frac2{\sqrt{2\pi}}\int_0^\infty t^se^{-t^2/2}\dd t=\frac2{\sqrt{2\pi}}\int_0^\infty(2u)^{(s-1)/2}e^{-u}\dd u$, which is the claim; for $s=2k$ use $\Gamma(\frac12)=\sqrt\pi$ and $\Gamma(x+1)=x\Gamma(x)$ (\cref{def:toolkit:gamma}). The moments of $\abs Y$ are finite because $\abs Y^k\le d^{k/2}\sum_i\abs{Y_i}^k$, and an odd polynomial $P$ has $\E P(Y)=\E P(-Y)=-\E P(Y)$ because $Y$ and $-Y$ have the same law.
(iii) Let $\phi(x)=(2\pi)^{-1/2}e^{-x^2/2}$, so that $\phi'(x)=-x\phi(x)$. Since $x/t\ge1$ on $x\ge t$, $\Prob(Z>t)\le\int_t^\infty\frac xt\phi(x)\dd x=\phi(t)/t$. Since $\phi(t+u)=\phi(u)e^{-tu-t^2/2}\le\phi(u)e^{-t^2/2}$ for $u\ge0$, $\Prob(Z>t)\le\frac12e^{-t^2/2}$. Since the derivative of $-\phi(x)/x$ is $(1+x^{-2})\phi(x)$, $\phi(t)/t=\int_t^\infty(1+x^{-2})\phi(x)\dd x\le(1+t^{-2})\Prob(Z>t)$.
(iv) By (i) with $A=\frac12I_d$, $\E e^{\abs Y^2/4}=2^{d/2}$, and Markov's inequality gives $\Prob(\abs Y\ge r)\le e^{-r^2/4}\E e^{\abs Y^2/4}$.
\end{proof}

\begin{proof}[Proof of \cref{thm:toolkit:laplace}]
The plan is to cut $W$ into a small ball around $w_0$ and its complement. Outside the ball, hypothesis (iii) makes the integral exponentially small. Inside, the substitution $w=w_0+u/\sqrt n$ and Taylor's theorem turn the integrand into a Gaussian times $g(w_0)$, plus a term of order $n^{-1/2}$ that is odd in $u$ and integrates to zero, plus a remainder of order $n^{-1}$. We may assume that $f(w_0)=0$. The degenerate case $g(w_0)=0$ is included; the theorem then says that the integral is $O(n^{-d/2-1})$.

Let $\mu>0$ be the smallest eigenvalue of $H$. On $\bar B(w_0,r_0)$ let $G$ bound $\abs g$, let $C_1$ and $C_2$ bound the norms of the first and second derivatives of $g$, and let $C_3$ and $C_4$ bound those of the third and fourth derivatives of $f$. Taylor's theorem gives $\abs{f(w_0+v)-\frac12v^\top Hv}\le\frac{C_3}6\abs v^3$ for $\abs v\le r_0$; with $\delta=\min(r_0,3\mu/(2C_3))$, or $\delta=r_0$ if $C_3=0$,
\begin{equation}\label{eq:toolkit:laplaceball}
\bigl\lvert f(w_0+v)-\tfrac12v^\top Hv\bigr\rvert\le\tfrac\mu4\abs v^2\qquad(\abs v\le\delta).
\end{equation}

\emph{Outside the ball.} For $n\ge n_0$ and $\abs{w-w_0}\ge\delta$, (iii) gives $e^{-nf(w)}\le e^{-(n-n_0)\eta(\delta)}e^{-n_0f(w)}$, so by (iv) the integral over $W\setminus B(w_0,\delta)$ is $O(e^{-n\eta(\delta)})$.

\emph{Inside the ball.} Substitute $w=w_0+u/\sqrt n$, and write $\rho_n(u)=nf(w_0+u/\sqrt n)-\frac12u^\top Hu$ and $g_n(u)=g(w_0+u/\sqrt n)$. The integral over the ball becomes $n^{-d/2}\int_{\abs u<\delta\sqrt n}g_n(u)e^{-\frac12u^\top Hu-\rho_n(u)}\dd u$. Taylor's theorem gives $\rho_n=n^{-1/2}T+r_n$, with the cubic form $T(u)=\frac16\sum_{ijk}\partial_{ijk}f(w_0)u_iu_ju_k$ and $\abs{r_n(u)}\le C_4\abs u^4/(24n)$, and $g_n=g(w_0)+n^{-1/2}\nabla g(w_0)\cdot u+s_n$ with $\abs{s_n(u)}\le C_2\abs u^2/(2n)$; also $\abs{g_n(u)-g(w_0)}\le C_1\abs u/\sqrt n$, and $\abs{\rho_n(u)}\le\frac\mu4\abs u^2$ by \eqref{eq:toolkit:laplaceball}. Writing $e^{-\rho}=1-\rho+E(\rho)$ with $\abs{E(\rho)}\le\frac12\rho^2e^{\abs\rho}$ and collecting terms,
\[
g_ne^{-\rho_n}=g(w_0)+n^{-1/2}P+R_n,\qquad P(u)=\nabla g(w_0)\cdot u-g(w_0)T(u),
\]
with $R_n=s_n-g(w_0)r_n-(g_n-g(w_0))\rho_n+g_nE(\rho_n)$. For $n\ge1$ and $\abs u<\delta\sqrt n$ these bounds give $\abs{R_n(u)}\le\frac cn(1+\abs u^8)e^{\mu\abs u^2/4}$, with $c$ independent of $n$ and $u$; the largest term is $g_nE(\rho_n)$, in which $\rho_n^2\le2(C_3^2\abs u^6/36+C_4^2\abs u^8/576)/n$. Since $e^{-\frac12u^\top Hu}\le e^{-\mu\abs u^2/2}$, the remainder contributes $O(n^{-1})$, by part (ii) of \cref{lem:toolkit:gaussian}. The odd polynomial $P$ integrates to zero against the even weight over the ball. By part (i), $\int_{\abs u<\delta\sqrt n}e^{-\frac12u^\top Hu}\dd u$ is $(2\pi)^{d/2}(\det H)^{-1/2}$ minus at most $\int_{\abs u\ge\delta\sqrt n}e^{-\mu\abs u^2/2}\dd u$, which part (iv) bounds by $(4\pi/\mu)^{d/2}e^{-n\mu\delta^2/4}$ after the substitution $u=y/\sqrt\mu$.

\emph{Collecting.} The integral over $W$ is $n^{-d/2}\bigl((2\pi)^{d/2}(\det H)^{-1/2}g(w_0)+O(n^{-1})\bigr)$ plus exponentially small terms. This is \eqref{eq:toolkit:laplace}, and taking logarithms gives the second display, because $\log(1+O(n^{-1}))=O(n^{-1})$.
\end{proof}

\Cref{ch:regular} proves a random version of \cref{thm:toolkit:laplace}, with $f=L_n$, for the free energy of a regular model (\cref{thm:regular:laplace}), and \cref{ch:coefficient} uses the theorem itself in an exercise to integrate out a noise variance; \cref{lem:toolkit:gaussian} serves wherever a Gaussian is integrated, as in \cref{ch:regular,ch:coefficient,ch:freeenergy}. At a singular point the proof breaks at its first step, because no ball satisfies \eqref{eq:toolkit:laplaceball}; the rest of this appendix collects what replaces it.

\begin{theorem}[Stein's identity and Gaussian smoothing]\label{thm:toolkit:stein}
Let $Z\sim\Normal(0,1)$.
\begin{enumerate}
\item (Stein's identity.) If $f\colon\R\to\R$ is absolutely continuous on every bounded interval and $\E\abs{f'(Z)}<\infty$, then $\E\abs{Zf(Z)}<\infty$ and $\E[Zf(Z)]=\E[f'(Z)]$. If $Y\sim\Normal(0,I_d)$ and $f\colon\R^d\to\R$ is continuously differentiable with $\E\abs{Y_if(Y)}$ and $\E\abs{\partial_if(Y)}$ finite, then $\E[Y_if(Y)]=\E[\partial_if(Y)]$.
\item (The heat equation.) Let $\psi\colon\R\to\R$ be measurable with $\int\abs{\psi(t)}e^{-\kappa t^2}\dd t<\infty$ for some $\kappa\in(0,\frac12)$, and let $\Psi(s,z)=\E\,\psi(z+\sqrt sZ)$ for $0<s<1/(2\kappa)$ and $z\in\R$. Then $\Psi$ is infinitely differentiable, its derivatives may be taken under the integral sign, and $\partial_s\Psi=\frac12\partial_z^2\Psi$.
\item (Tweedie's formula.) If moreover $\psi\ge0$ and $\psi>0$ on a set of positive measure, then $h(z)=\int e^{-(t-z)^2/2}\psi(t)\dd t$ is positive, and the probability density proportional to $e^{-(t-z)^2/2}\psi(t)$ has mean $z+h'(z)/h(z)$ and variance $1+(\log h)''(z)$.
\item (Homogeneity.) For $\alpha>-1$, the function $h(z)=\int e^{-(t-z)^2/2}\abs t^\alpha\dd t$ satisfies
\begin{equation}\label{eq:toolkit:homogeneity}
h''(z)=\alpha h(z)-zh'(z)\quad(z\in\R),\qquad h(0)=2^{(\alpha+1)/2}\,\Gamma\bigl(\tfrac{\alpha+1}2\bigr),
\end{equation}
and $r=h'/h$ satisfies $r'=\alpha-zr-r^2$.
\end{enumerate}
\end{theorem}

\begin{proof}
The plan is to prove (i) by exchanging two integrals and (ii) by differentiating the heat kernel, and to deduce (iii) and (iv) from (ii). (i) Let $\phi$ be the standard normal density; since $\phi'(x)=-x\phi(x)$, $\phi(z)=\int_z^\infty x\phi(x)\dd x=-\int_{-\infty}^zx\phi(x)\dd x$. Hence
\[
\E f'(Z)=\int_0^\infty f'(z)\int_z^\infty x\phi(x)\dd x\dd z-\int_{-\infty}^0f'(z)\int_{-\infty}^zx\phi(x)\dd x\dd z.
\]
With $\abs{f'}$ in place of $f'$ the right side is $\E\abs{f'(Z)}<\infty$, so Fubini's theorem exchanges the integrals, and $\int_0^xf'=f(x)-f(0)$ turns the right side into $\E[Z(f(Z)-f(0))]=\E[Zf(Z)]$; the same computation bounds $\E\abs{Z(f(Z)-f(0))}$ by $\E\abs{f'(Z)}$. In the vector case, for almost every value of the coordinates other than $Y_i$ the one-variable identity holds in $Y_i$ (Fubini), and integrating over those coordinates gives the claim.
(ii) Write $\Psi(s,z)=\int p_s(z-t)\psi(t)\dd t$ with $p_s(x)=(2\pi s)^{-1/2}e^{-x^2/(2s)}$, for which $\partial_sp_s=(\frac{x^2}{2s^2}-\frac1{2s})p_s=\frac12\partial_x^2p_s$. Every derivative of $p_s(z-t)$ in $(s,z)$ is a polynomial in $z-t$, $s^{-1/2}$ and $s^{-1}$ times $p_s(z-t)$. For $s_1\le s\le s_2<1/(2\kappa)$ and $\abs z\le A$, the inequality $(t-z)^2\ge(1-\theta)t^2-(\theta^{-1}-1)A^2$, with $\theta>0$ so small that $(1-\theta)/(2s_2)>\kappa$, bounds each such derivative by a constant times $e^{-\kappa t^2}$; so dominated convergence allows differentiation under the integral sign, and the heat equation passes from $p_s$ to $\Psi$.
(iii) Here $h(z)=\sqrt{2\pi}\,\Psi(1,z)$, which is positive. By (ii), $h'(z)=\int(t-z)e^{-(t-z)^2/2}\psi(t)\dd t$ and $h''(z)=\int((t-z)^2-1)e^{-(t-z)^2/2}\psi(t)\dd t$, so under the density of (iii) the mean of $t-z$ is $h'/h$ and the mean of $(t-z)^2$ is $1+h''/h$; the variance is $1+h''/h-(h'/h)^2=1+(\log h)''$.
(iv) The function $\abs t^\alpha$ is locally integrable because $\alpha>-1$, so (ii) applies, and the substitution $t=\sqrt s\,\tau$ gives the homogeneity $\Psi(s,z)=(2\pi)^{-1/2}s^{\alpha/2}h(z/\sqrt s)$. At $s=1$, $\partial_s\Psi=(2\pi)^{-1/2}(\frac\alpha2h-\frac z2h')$ and $\partial_z^2\Psi=(2\pi)^{-1/2}h''$, so the heat equation reads $h''=\alpha h-zh'$. Finally $h(0)=\sqrt{2\pi}\,\E\abs Z^\alpha$, which part (ii) of \cref{lem:toolkit:gaussian} evaluates, and $r'=h''/h-r^2=\alpha-zr-r^2$.
\end{proof}

Identity \eqref{eq:toolkit:homogeneity} can also be checked directly, since $\alpha h-zh'-h''$ is the integral of the derivative of $t\abs t^\alpha e^{-(t-z)^2/2}$. That integration by parts in $t$ is the virial identity of \cref{ch:generalization}, which proves the analogue of \eqref{eq:toolkit:homogeneity} for every inverse temperature and for unequal weights on the two sides of $0$ (\cref{lem:generalization:parts}). There, for a one-coefficient model whose coefficient density behaves like $\abs c^\alpha$ at the truth, the mean of the rescaled limit posterior is $M(Z)=Z+r(Z)$ with $r=h'/h$, as in (iii), and the virial identity together with Stein's identity (i), applied to $M$, gives $\E[M(Z)^2]=1+\alpha$, which is $2\lambda$ when $\alpha=2\lambda-1$. With (iv) the same value follows from (i) applied to $r$: $\E[Zr]=\E[r']=\alpha-\E[Zr]-\E[r^2]$, so $\E[(Z+r)^2]=1+2\E[Zr]+\E[r^2]=1+\alpha$ (\cref{ex:generalization:riccati}). We confirmed \eqref{eq:toolkit:homogeneity} and this value by quadrature for seven values of $\alpha$ between $-\frac34$ and $2$.\sidenote{The identity underlies \textcite{Stein1981}'s unbiased estimate of risk; \textcite{Efron2011} discusses Tweedie's formula and its history.}\index{Stein's identity}\index{Tweedie's formula}

\section{The Gamma function, Mellin transforms, and regular variation}\label{sec:toolkit:transforms}

The learning coefficient enters the theory through two transforms of the law of $K(w)$ under the prior: the zeta function $\zeta(z)=\int K^z\varphi$, a Mellin transform, whose largest pole is $-\lambda$; and the partition function $\int e^{-nK}\varphi$, a Laplace transform, which behaves like $n^{-\lambda}(\log n)^{m-1}$.

\begin{definition}[The Gamma function]\label{def:toolkit:gamma}
For complex $s$ with $\Re s>0$, the \term{Gamma function}\index{Gamma function} is $\Gamma(s)=\int_0^\infty t^{s-1}e^{-t}\dd t$. It has the following properties.
\begin{enumerate}
\item $\Gamma(s+1)=s\Gamma(s)$ and $\Gamma(1)=1$, so that $\Gamma(k+1)=k!$ for integers $k\ge0$; and $\Gamma(\frac12)=\sqrt\pi$, $\Gamma(\frac32)=\frac{\sqrt\pi}2$.
\item For $n>0$ and $\Re s>0$, $\int_0^\infty t^{s-1}e^{-nt}\dd t=\Gamma(s)\,n^{-s}$.
\item For $\Re s>-1$ and integers $k\ge0$, $\int_0^1t^s(\log\frac1t)^k\dd t=k!/(s+1)^{k+1}$.
\end{enumerate}
\end{definition}

\begin{proof}[Proofs]
(i) Integrating by parts, $\Gamma(s+1)=[-t^se^{-t}]_0^\infty+s\int_0^\infty t^{s-1}e^{-t}\dd t=s\Gamma(s)$, and $\Gamma(1)=\int_0^\infty e^{-t}\dd t=1$; the substitution $t=u^2$ gives $\Gamma(\frac12)=2\int_0^\infty e^{-u^2}\dd u=\sqrt\pi$ by part (i) of \cref{lem:toolkit:gaussian}. (ii) Substitute $t'=nt$. (iii) The substitution $t=e^{-u}$ turns the integral into $\int_0^\infty u^ke^{-(s+1)u}\dd u$, which is $k!/(s+1)^{k+1}$ by (ii) for real $s>-1$; both sides are holomorphic in $s$ on $\Re s>-1$, so they agree there.
\end{proof}

Integrals of this kind recur throughout the book: part (ii) gives the constant $\Gamma(\lambda+1)$ of \cref{thm:toolkit:karamata}, and part (iii) the Mellin transforms below.

\begin{definition}[Mellin transform]\label{def:toolkit:mellin}
The \term{Mellin transform}\index{Mellin transform} of a measurable function $v$ on $(0,\infty)$ is $\cM v(z)=\int_0^\infty t^zv(t)\dd t$, at the complex $z$ for which the integral converges absolutely; for a finite measure $\nu$ on $(0,\infty)$ it is $\int t^z\nu(\dd t)$. It has the following properties.
\begin{enumerate}
\item For real $s>-1$ and integers $k\ge0$, the transform of $t^s(\log\frac1t)^k\ind_{(0,1]}(t)$ is $k!/(z+s+1)^{k+1}$ on $\Re z>-s-1$; it continues to a meromorphic function on $\C$ with one pole, of order $k+1$, at $z=-s-1$.
\item (Continuation by integrating by parts.) Let $N\ge1$ and $k\ge0$ be integers and let $b$ have $N$ continuous derivatives on $[0,1]$. Then $I_k(z)=\int_0^1t^z(\log\frac1t)^kb(t)\dd t$, which is holomorphic on $\Re z>-1$, continues to a meromorphic function on $\Re z>-N-1$ whose poles lie among $-1,\dots,-N$; near $z=-(j+1)$ it equals $k!\,b^{(j)}(0)/\bigl(j!\,(z+j+1)^{k+1}\bigr)$ plus a holomorphic function. If $b$ is infinitely differentiable, then $I_k$ is meromorphic on $\C$.
\item (Uniqueness.) If two functions integrable on $(0,1]$ and zero on $(1,\infty)$ have transforms that agree on a vertical line $\Re z=\sigma\ge0$, then the functions agree almost everywhere.
\end{enumerate}
\end{definition}

The convention is shifted by one from the usual $\int t^{s-1}v(t)\dd t$, so that the zeta function, $\int K^z\varphi=\int t^z\nu(\dd t)$ with $\nu$ the law of $K(w)$ under $\varphi$, is the Mellin transform of $\nu$, and of the state density $v=V'$ when it exists. Part (ii) is the mechanism behind every pole in the book: a function that behaves like $t^s(\log\frac1t)^k$ near $0$ has its rightmost pole at $-(s+1)$, of order $k+1$, and nothing away from $0$ contributes a pole; for a state density $v(t)\approx Ct^{\lambda-1}(\log\frac1t)^{m-1}$ this is the pole of order $m$ at $-\lambda$. Transforms also turn products into products: for independent positive random variables, $\E[(XX')^z]=\E[X^z]\,\E[X'^z]$. For two independent uniforms on $(0,1]$ the product therefore has transform $(z+1)^{-2}$, which by (i) and (iii) is the transform of the density $\log\frac1t$; in the same way the product of $L$ of them has the density $(\log\frac1t)^{L-1}/(L-1)!$, which \cref{ch:coefficient} finds by the equivalent route that $-\log\abs{a_1\cdots a_L}$ is a sum of $L$ independent exponential variables.

\begin{proof}[Proof of (ii) and (iii)]
The plan is to integrate by parts once, which moves the half-plane of convergence one unit to the left, to iterate by induction on $N$, and to obtain the powers of the logarithm by differentiating in $z$. Write $I_0(z;b)=\int_0^1t^zb(t)\dd t$. For $\Re z>-1$, integrating $t^z=\frac{\dd}{\dd t}\frac{t^{z+1}}{z+1}$ by parts against $b$ gives
\begin{equation}\label{eq:toolkit:mellinparts}
I_0(z;b)=\frac{b(1)}{z+1}-\frac{I_0(z+1;b')}{z+1};
\end{equation}
the boundary term at $0$ vanishes because $\Re(z+1)>0$. Since $b'$ is bounded, $I_0(z+1;b')$ is holomorphic on $\Re z>-2$, so \eqref{eq:toolkit:mellinparts} continues $I_0(\cdot;b)$ to $\Re z>-2$ with at most a simple pole at $-1$, of residue $b(1)-I_0(0;b')=b(0)$. This is the case $N=1$, $k=0$. Suppose the case $k=0$ holds for $N-1$, and let $b$ have $N$ continuous derivatives. By the induction hypothesis applied to $b'$, the function $z\mapsto I_0(z+1;b')$ is meromorphic on $\Re z>-N-1$, with poles among $-2,\dots,-N$ and residue $b^{(j)}(0)/(j-1)!$ at $-(j+1)$. By \eqref{eq:toolkit:mellinparts}, $I_0(\cdot;b)$ is meromorphic there too, with residue $-\frac1{-j}\cdot\frac{b^{(j)}(0)}{(j-1)!}=\frac{b^{(j)}(0)}{j!}$ at $-(j+1)$. For $k\ge1$, since $(\log\frac1t)^kt^z=(-\partial_z)^kt^z$ and differentiation under the integral sign is allowed on $\Re z>-1$, $I_k=(-1)^kI_0^{(k)}$; and $(-\partial_z)^k$ maps $c/(z+j+1)$ to $k!\,c/(z+j+1)^{k+1}$, and holomorphic functions to holomorphic functions. This proves (ii). For (iii), the substitution $t=e^{-u}$ turns the transform on the line $\Re z=\sigma$ into the Fourier transform of the integrable function $u\mapsto e^{-(\sigma+1)u}v(e^{-u})$ on $(0,\infty)$, which determines it almost everywhere.
\end{proof}

Applied in each coordinate in turn, with the other coordinates as parameters, part (ii) shows that $\int_{[0,1]^d}u_1^{2k_1z+h_1}\cdots u_d^{2k_dz+h_d}\,b(u)\dd u$, for $b$ infinitely differentiable, is meromorphic on $\C$ with poles among the numbers $-(h_j+1+\ell)/(2k_j)$, $\ell=0,1,2,\dots$; with \cref{thm:toolkit:hironaka} this is why the poles of the zeta function are negative rational numbers. \Cref{ch:coefficient} meets the Mellin transform as the zeta function of a one-coefficient model (\cref{prop:coefficient:zeta}), \cref{ch:volume} uses it for the relation between the zeta function and the volume (\cref{thm:volume:zeta}), and \cref{ch:blowups,ch:newton} for the zeta function on the charts of a resolution.

A positive function $\ell$ on $(0,\infty)$ is \term{slowly varying} at infinity if $\ell(cx)/\ell(x)\to1$ as $x\to\infty$ for every $c>0$, as powers of $\log x$ are. A function $V$ on $(0,\infty)$ is \term{regularly varying}\index{regular variation} at $0$ with index $\lambda$ if $V(\eps)=\eps^\lambda\ell(1/\eps)$ with $\ell$ slowly varying.

\begin{theorem}[Karamata's theorem for Laplace transforms]\label{thm:toolkit:karamata}
Let $\nu$ be a finite measure on $[0,\infty)$ with $\nu(\{0\})=0$, let $V(\eps)=\nu([0,\eps))$ for $\eps>0$, and let $Z(n)=\int e^{-nt}\nu(\dd t)$ for $n>0$.
\begin{enumerate}
\item (Abelian direction, the case of this book.) If $V(\eps)\sim C\eps^\lambda(\log\frac1\eps)^{m-1}$ as $\eps\to0$, with constants $C>0$, $\lambda>0$ and $m\ge1$, then as $n\to\infty$
\begin{gather*}
Z(n)\sim C\,\Gamma(\lambda+1)\,n^{-\lambda}(\log n)^{m-1},\\
\int te^{-nt}\nu(\dd t)\sim C\lambda\,\Gamma(\lambda+1)\,n^{-\lambda-1}(\log n)^{m-1}.
\end{gather*}
Consequently $-\log Z(n)=\lambda\log n-(m-1)\log\log n-\log\bigl(C\Gamma(\lambda+1)\bigr)+o(1)$, and the mean of $nt$ under the probability measure proportional to $e^{-nt}\nu(\dd t)$ tends to $\lambda$.
\item (Abelian direction, general.) If $V$ is regularly varying at $0$ with index $\lambda\ge0$, then $Z(n)\sim\Gamma(\lambda+1)\,V(1/n)$.
\item (Tauberian direction.) Conversely, if $Z(n)\sim\Gamma(\lambda+1)\,n^{-\lambda}\ell(n)$ as $n\to\infty$, with $\lambda\ge0$ and $\ell$ slowly varying at infinity, then $V(\eps)\sim\eps^\lambda\ell(1/\eps)$ as $\eps\to0$. In particular $Z(n)\sim C'n^{-\lambda}(\log n)^{m-1}$ implies $V(\eps)\sim\frac{C'}{\Gamma(\lambda+1)}\eps^\lambda(\log\frac1\eps)^{m-1}$.
\end{enumerate}
\end{theorem}

When $\nu$ is the law of $K(w)$ under the prior, $V$ is the book's volume function and $Z(n)=\int e^{-nK}\varphi$, so that $-\log Z(n)$ is the population free energy $\bar F(n)$.

\begin{proof}[Proof of (i)]
The plan is to write both transforms as integrals of $V$ against explicit kernels, to rescale so that $V$ is evaluated at $v/n$, and to pass to the limit by dominated convergence. For $t\ge0$, $e^{-nt}=\int_t^\infty ne^{-ns}\dd s$, so Tonelli's theorem gives $Z(n)=n\int_0^\infty e^{-ns}\nu([0,s))\dd s=n\int_0^\infty e^{-ns}V(s)\dd s$. In the same way $te^{-nt}=\int_t^\infty(ns-1)e^{-ns}\dd s$, and Fubini's theorem applies because $\nu$ is finite; hence $\int te^{-nt}\nu(\dd t)=\int_0^\infty(ns-1)e^{-ns}V(s)\dd s$. With $J_k(n)=\int_0^\infty v^ke^{-v}V(v/n)\dd v$, the substitution $s=v/n$ gives
\[
Z(n)=J_0(n),\qquad\int te^{-nt}\nu(\dd t)=\frac{J_1(n)-J_0(n)}n.
\]
It therefore suffices to show that $J_k(n)\sim C\Gamma(\lambda+k+1)\,a_n$ for $k=0,1$, with $a_n=n^{-\lambda}(\log n)^{m-1}$; then $J_1-J_0\sim C(\Gamma(\lambda+2)-\Gamma(\lambda+1))a_n=C\lambda\Gamma(\lambda+1)a_n$, and the ratio of the two transforms gives the mean of $nt$.

Choose $\eps_0\in(0,1)$ with $V(\eps)\le2C\eps^\lambda(\log\frac1\eps)^{m-1}$ for $0<\eps<\eps_0$, and let $n\ge e$. On $v\ge n\eps_0$ we use only $V\le M=\nu([0,\infty))$: that part of $J_k(n)$ is at most $M\int_{n\eps_0}^\infty v^ke^{-v}\dd v$, which decays exponentially and is $o(a_n)$. On $0<v<n\eps_0$, the integrand of $J_k(n)/a_n$ converges to $Cv^{k+\lambda}e^{-v}$ for each $v$, because $V(v/n)\sim C(v/n)^\lambda(\log\frac nv)^{m-1}$ and $\log\frac nv/\log n\to1$. It is also dominated: since $\log\frac nv/\log n=1+\log\frac1v/\log n$ is at most $1$ for $v\ge1$ and at most $1+\log\frac1v$ for $v<1$, it is at most
\[
2C\,v^{k+\lambda}e^{-v}\bigl(1+\log^+\tfrac1v\bigr)^{m-1},
\]
which is integrable on $(0,\infty)$ because $\lambda>0$. Dominated convergence gives $J_k(n)/a_n\to C\int_0^\infty v^{k+\lambda}e^{-v}\dd v=C\Gamma(\lambda+k+1)$.
\end{proof}

Parts (ii) and (iii) are Karamata's theorem, which we state without proof. \Textcite[Theorem~1.7.1]{BinghamEtAl1987} state it for a nondecreasing function at infinity and its Laplace--Stieltjes transform near $0$, as does \textcite[Chapter~XIII]{Feller1971}. The version at the other end, used here, has the same proof: the measures $\nu_\tau(A)=\nu(\tau^{-1}A)/Z(\tau)$ have the transforms $Z(s\tau)/Z(\tau)$, which tend to $s^{-\lambda}$ as $\tau\to\infty$ under the hypothesis of (iii), and the continuity theorem for Laplace transforms turns this into $\nu_\tau([0,t])\to t^\lambda/\Gamma(\lambda+1)$, which at $t=1$ is (iii). Part (ii) is proved like part (i), with Potter's bounds for slowly varying functions \parencite{BinghamEtAl1987} in place of the explicit form of $V$ in the domination. For the zeta function, \cref{ch:volume} proves a Tauberian statement of another kind, from the largest pole of $\zeta$ to the asymptotics of $V$, with a theorem of \textcite{Delange1954} (\cref{thm:volume:zeta}).\index{Karamata's theorem} The limit in (i) is approached slowly when $m\ge2$: for the product model, $K=a^2b^2/2$ under the uniform prior on $[-1,1]^2$, $C=1/\sqrt2$, $\lambda=\frac12$, $m=2$, and $Z(n)=\frac12\sqrt{\pi/(2n)}\,(\log n+\gamma_{\mathrm E}+\log2)$ up to an exponentially small error, with $\gamma_{\mathrm E}$ Euler's constant, so the ratio of $Z(n)$ to its limit form is still $1.07$ at $n=10^8$. Part (i) is \cref{lem:volume:abelian} of \cref{ch:volume}, and \cref{ch:freeenergy} applies it to multiples of $K$; its second asymptotic, which says that the tempered posterior of the population loss holds the energy $\E^\beta_w[nK]\approx\lambda/\beta$, is refined in \cref{prop:temperature:mean} and underlies the estimators of \cref{ch:temperature,ch:llc}.

\section{Analytic functions and the resolution of singularities}\label{sec:toolkit:singularities}

Watanabe's theory needs the loss to be analytic, because analytic functions are the functions whose singularities can be resolved. The splitting lemma at the end of the section separates the regular directions of a critical point from the singular ones.

\begin{definition}[Real analytic functions]\label{def:toolkit:analytic}
Let $U\subset\R^d$ be open. A function $f\colon U\to\R$ is \term{real analytic}\index{real analytic function} if every point $w^*\in U$ has a neighborhood on which $f$ is the sum of a convergent power series $\sum_\alpha c_\alpha(w-w^*)^\alpha$, where $\alpha$ runs over the multi-indices in $\N^d$ and $(w-w^*)^\alpha=\prod_j(w_j-w^*_j)^{\alpha_j}$. A function on a set $W\subset\R^d$ that is not open, such as a compact parameter set, is analytic if it is the restriction of an analytic function on an open set containing $W$; a map into $\R^k$ is analytic if its coordinates are. An \term{analytic set} is a set $\{w\in U:f_1(w)=\dots=f_k(w)=0\}$ with analytic $f_j$; when $K$ is analytic, the set of true parameters $W_0=\{K=0\}$ is one. A $d$-dimensional real analytic manifold is a space covered by charts onto open subsets of $\R^d$ with analytic transition maps; an analytic isomorphism is an analytic bijection with an analytic inverse; and a continuous map is proper if the preimage of every compact set is compact. The book uses three facts.
\begin{enumerate}
\item Analytic functions are infinitely differentiable, and their power series are their Taylor series. Sums, products, compositions, and quotients with nonvanishing denominators of analytic functions are analytic; polynomials, $\exp$, $\tanh$, and $\log$ on $(0,\infty)$ are analytic.
\item (Identity theorem.) If $U$ is connected and $f$ is analytic on $U$ and vanishes on a nonempty open subset, then $f=0$ on $U$.
\item If $U$ is connected and $f$ is analytic and not identically zero on $U$, then $\{f=0\}$ has Lebesgue measure zero.
\end{enumerate}
\end{definition}

\begin{proof}[Proof of (ii) and (iii)]
The plan is to prove (ii) by a connectedness argument and to reduce (iii), through (ii), to the zeros of functions of one variable. (ii) The set of points at which every derivative of $f$ vanishes is closed, because the derivatives are continuous; it is open, because at such a point the Taylor series is zero, so $f$ and its derivatives vanish on a neighborhood; and it is nonempty, because it contains the open set on which $f$ vanishes. Since $U$ is connected, it is all of $U$.
(iii) Because every open cover of $U$ has a countable subcover, it suffices to show that each point $w^*$ has a neighborhood in which $\{f=0\}$ is null. By (ii), some derivative of $f$ is nonzero at $w^*$; let $k$ be the smallest order of such a derivative. Then $P(v)=\sum_{\abs\alpha=k}\partial^\alpha f(w^*)v^\alpha/\alpha!$ is a nonzero polynomial, so some unit vector $v$ has $P(v)\ne0$, and the $k$-th derivative of $f$ at $w^*$ in the direction $v$, which is $k!\,P(v)$, is not zero. After a rotation, $v$ is the first coordinate vector, and $\partial_1^kf\ne0$ on a small cube $Q$ around $w^*$. For fixed $(w_2,\dots,w_d)$, the function $w_1\mapsto f(w)$ on the side of $Q$ has a nonvanishing $k$-th derivative, so by Rolle's theorem it has at most $k$ zeros; by Fubini's theorem, $\{f=0\}\cap Q$ is null.
\end{proof}

By (iii), the set of true parameters of an analytic $K$ that is not identically zero is null, so $V(\eps)\to0$ as $\eps\to0$ for every prior with a density. By (ii), an analytic function on $\R$ that vanishes on $(-1,0)$ vanishes everywhere, so the ReLU is not analytic; the ReLU neuron of the running models escapes the problem because its $K=a^2b^2/4$ is a polynomial. Analyticity can also fail more quietly. For the tanh neuron with $x\sim\Normal(0,1)$, $F(b)=\E[\tanh^2(bx)]$ is infinitely differentiable, but its Taylor coefficients at $0$ are those of $\tanh^2$, of order $k(2/\pi)^{2k}$, times $\E x^{2k}=(2k-1)(2k-3)\cdots1$, so its Taylor series diverges for every $b\ne0$ and $K=\frac{a^2}2F(b)$ is not analytic at the singular point (with inputs bounded by $R$ it is analytic for $\abs b<\pi/(2R)$). Its learning coefficient is still that of $a^2b^2$, by comparison, but theorems that assume analyticity do not apply to it directly. \Cref{ch:singular,ch:volume,ch:blowups,ch:llc} use these facts.\index{analyticity}

\begin{theorem}[Hironaka's resolution of singularities, in the form used by Watanabe]\label{thm:toolkit:hironaka}
Let $f$ be a real analytic function on a neighborhood of a point $w^*\in\R^d$, with $f(w^*)=0$ and $f$ not identically zero near $w^*$. Then there are an open neighborhood $U$ of $w^*$, a $d$-dimensional real analytic manifold $M$, and a proper real analytic map $g\colon M\to U$ with the following properties.
\begin{enumerate}
\item Let $U_0=\{w\in U:f(w)=0\}$ and $M_0=g^{-1}(U_0)$. Then $g$ restricted to $M\setminus M_0$ is an analytic isomorphism onto $U\setminus U_0$.
\item Every point $P\in M_0$ has analytic local coordinates $u=(u_1,\dots,u_d)$ centered at $P$ in which
\[
f(g(u))=S\,u_1^{k_1}\cdots u_d^{k_d},\qquad\det g'(u)=b(u)\,u_1^{h_1}\cdots u_d^{h_d},
\]
with a sign $S\in\{-1,1\}$, nonnegative integers $k_j$ and $h_j$, not all $k_j$ zero, and an analytic function $b$ with $b(0)\ne0$.
\end{enumerate}
\end{theorem}

The theorem is due to \textcite{Hironaka1964a,Hironaka1964b}, who proved resolution of singularities for algebraic varieties over fields of characteristic zero. The real analytic local form above, with the Jacobian, is the version that \textcite[Theorem~3.3]{Lin2011} attributes to \textcite{Atiyah1970} and that, by Lin's account, \textcite{Watanabe2009} uses as his Theorem~2.3; we have not checked either attribution in the original. \Textcite[Section~3]{Watanabe2024} states the theorem directly for a loss and a prior, in the form \eqref{eq:toolkit:standard} below. The book uses it as a black box. A function is \term{normal crossing}\index{normal crossing} at a point $P$ if, in some analytic coordinates centered at $P$, it equals a monomial $u_1^{k_1}\cdots u_d^{k_d}$ times an analytic function that does not vanish at $P$; near $P$ its zero set is then the union of the coordinate hyperplanes $\{u_j=0\}$ with $k_j>0$, which cross each other transversally; the theorem makes every analytic function normal crossing, together with the Jacobian, by a change of variables that changes nothing away from the zero set. The product model's $K=a^2b^2/2$ is normal crossing at the origin already; $a^2+b^4$ is not, because its zero set near the origin is a single point, while the zero set of a normal-crossing function is a union of hypersurfaces; \cref{ch:blowups} resolves it by blow-ups.\sidenote{Hironaka received the Fields Medal for this work in 1970 \parencite{IMU1970}.}

For a loss $K\ge0$ and a prior $\varphi=\varphi_1\varphi_2$, with $\varphi_1\ge0$ analytic and $\varphi_2>0$ infinitely differentiable, the theorem takes the form in which the book uses it. Applied to the product $K\varphi_1$, it makes $K$ and $\varphi_1$ normal crossing in the same coordinates, because a factor of a monomial times a unit is again a monomial times a unit. Since $K\circ g\ge0$, each exponent of $K\circ g$ is even and its sign is $+1$, for otherwise $K\circ g$ would change sign with one coordinate. With compactness this gives what \textcite[Theorem~1]{Watanabe2013} calls the \term{standard representation}\index{standard representation}: in each of finitely many coordinate cubes,
\begin{equation}\label{eq:toolkit:standard}
K(g(u))=u^{2k}=u_1^{2k_1}\cdots u_d^{2k_d},\qquad\varphi(g(u))\,\abs{\det g'(u)}=b(u)\,\abs{u^h},
\end{equation}
with $b>0$; that paper states that its proof is given in Theorem~6.1 of the 2009 book. The learning coefficient $\lambda$ is the minimum of $(h_j+1)/(2k_j)$ over all cubes and all $j$ with $k_j>0$, and the multiplicity $m$ is the largest number of indices $j$ attaining it in one cube (\cref{thm:blowups:pullback}). The zeta function becomes a finite sum of integrals of the type treated after \cref{def:toolkit:mellin}, so its poles are negative rational numbers; \textcite{BernsteinGelfand1969} and \textcite{Atiyah1970} used Hironaka's theorem in this way to continue $\int\abs f^z\varphi$ meromorphically, answering a question that I.\,M.\ Gel'fand had asked in 1954 (\cref{ch:volume} tells the story).\index{resolution of singularities}

\Cref{ch:blowups,ch:newton} carry out resolutions by hand and by toric changes of variables, \cref{ch:freeenergy,ch:generalization,ch:waic,ch:temperature} use the theorem through \cref{thm:toolkit:empirical}, and \cref{ch:llc} uses it to define the local learning coefficient.

\begin{theorem}[The splitting lemma]\label{thm:toolkit:splitting}
Let $f$ be real analytic near a point $w^*\in\R^d$ with $\nabla f(w^*)=0$, and let the Hessian $\nabla^2f(w^*)$ have rank $r$, with $p$ positive and $r-p$ negative eigenvalues. Then there are an analytic function $g$ near $0\in\R^{d-r}$ with $g(0)=0$, $\nabla g(0)=0$ and $\nabla^2g(0)=0$, and an analytic isomorphism $w=\Phi(x,y)$, $x\in\R^r$, $y\in\R^{d-r}$, from a neighborhood of $0$ onto a neighborhood of $w^*$ with $\Phi(0)=w^*$, such that
\[
f(\Phi(x,y))=f(w^*)+\sum_{i\le p}x_i^2-\sum_{p<i\le r}x_i^2+g(y).
\]
If $w^*$ is a local minimum of $f$, then $p=r$ and $g\ge0$ near $0$. For $r=d$ this is the Morse lemma: near a nondegenerate critical point, $f$ is a sum of signed squares in suitable coordinates.
\end{theorem}

\begin{proofsketch}
After a linear change of coordinates, $\nabla^2f(w^*)$ is diagonal with entries $\pm2$ in the first $r$ places and $0$ in the others; write $w-w^*=(x,y)$. Because $\partial_x^2f(w^*)$ is invertible, the analytic implicit function theorem solves $\partial_xf(x,y)=0$ by an analytic $x=\xi(y)$ with $\xi(0)=0$, and $\xi'(0)=0$ because the off-diagonal block of the Hessian vanishes. Taylor's formula with integral remainder in $x$ around $\xi(y)$ writes $f(\xi(y)+x',y)-f(\xi(y),y)=x'^\top Q(x',y)x'$ with $Q$ analytic and symmetric and $Q(0,0)=\diag(\pm1)$; completing the squares one variable at a time, which needs only square roots of analytic functions that do not vanish at $0$, turns this into $\sum_{i\le r}\pm x_i''^2$ by an analytic change $x''=T(x',y)x'$. The function $g(y)=f(\xi(y),y)-f(w^*)$ has vanishing first and second derivatives at $0$ because $\partial_xf(\xi(y),y)=0$, $\xi'(0)=0$ and the $y$-block of the Hessian vanishes. At a local minimum the Hessian has no negative eigenvalue, so $p=r$, and $g(y)=f(\Phi(0,y))-f(w^*)\ge0$. These are the steps of the usual proofs of the splitting lemma for smooth functions, which \textcite{GromollMeyer1969} prove on Hilbert spaces, and of the Morse lemma \parencite{ArnoldEtAl1985}; what we add is only that every step preserves analyticity. We omit two routine verifications: the induction that completes the squares with $y$ as a parameter, and the check that the composite of the three changes of variables has an invertible derivative at $0$, so that it is an analytic isomorphism near $0$ by the analytic inverse function theorem. The analytic implicit and inverse function theorems are in \textcite[Chapter~2]{KrantzParks2002}.
\end{proofsketch}

\Cref{ch:blowups,ch:newton} use the case $r=d$, the Morse lemma, for their regular answers. \Cref{ch:llc} proves its bounds $\frac r2\le\lambda(\wstar)\le\frac r2+\frac{d-r}4$ in terms of the rank $r$ of the Hessian (\cref{thm:llc:hessian}) by comparing volumes, for every loss $K\ge0$ with four continuous derivatives, without the splitting lemma and without analyticity; it uses the lemma only in \cref{cor:llc:minimal}, which says that for analytic $K$ the lower bound is attained exactly when the set of true parameters is, near $\wstar$, an analytic submanifold of dimension $d-r$.\index{splitting lemma}

\section{The log-likelihood ratio on a resolved space}\label{sec:toolkit:process}

Hironaka's theorem controls the deterministic part of the free energy. Watanabe's theorems also need the random part, the deviation of $nK_n$ from $nK$, and on the parameter space itself its natural normalization fails exactly where it matters. Write $f(x,w)=\log\{p_0(x)/p(x\mid w)\}$ with $p_0=p(\cdot\mid w_0)$, so that $K(w)=\E_X[f(X,w)]$ and $K_n(w)=\frac1n\sum_if(X_i,w)$ are the book's $K$ and $K_n$. Then
\[
nK_n(w)=nK(w)-\sqrt{nK(w)}\,\eta_n(w),\qquad\eta_n(w)=\frac1{\sqrt n}\sum_{i=1}^n\frac{K(w)-f(X_i,w)}{\sqrt{K(w)}},
\]
and $\eta_n$ is $0/0$ on the set of true parameters, where the posterior concentrates. After resolution, under the relatively finite variance of condition (iv) of \cref{ass:toolkit:fundamental} below, the numerator becomes divisible by the square root of the denominator, and the quotient is a well-defined random process with a Gaussian limit \parencite[Section~4.2]{Watanabe2024}.\sidenote{Watanabe presents the theory as the solution of these two problems: the singular geometry of $\{K=0\}$ and the $0/0$ in $\eta_n$ \parencite[Section~1]{Watanabe2024}.} The conditions are the following.

\begin{assumption}[Watanabe's fundamental conditions]\label{ass:toolkit:fundamental}
\leavevmode
\begin{enumerate}
\item The parameter set $W\subset\R^d$ is compact with nonempty interior, and $W=\{w:\pi_1(w)\ge0,\dots,\pi_k(w)\ge0\}$ for analytic functions $\pi_j$.
\item The prior is $\varphi=\varphi_1\varphi_2$ with $\varphi_1\ge0$ analytic and $\varphi_2>0$ infinitely differentiable.
\item Some $w_0$ in the interior of $W$ minimizes $L$, and all minimizers give the same density $p_0$. For some $s\ge6$, the map $w\mapsto f(\cdot,w)$ is analytic on an open set containing $W$ as a map into the space $L^s(q)$ of functions with $\int\abs{f(x)}^sq(x)\dd x<\infty$: near each point it is a power series in $w$ with coefficients in $L^s(q)$.
\item (Relatively finite variance.) There are constants $\eps_0,c>0$ such that $\E_X[f(X,w)]\ge c\,\E_X[f(X,w)^2]$ for every $w$ with $K(w)\le\eps_0$.
\end{enumerate}
\end{assumption}

These are the Fundamental Conditions (1)--(4) of \textcite[p.~872]{Watanabe2013}, with the interior minimizer and the uniqueness of $p_0$ that the paper assumes on p.~871. Condition (iv) holds when $q$ is realizable by the model or regular for it \parencite[p.~872]{Watanabe2013}; for an unrealizable and singular pair it can fail, and \textcite{Watanabe2010d} describes an example in which it fails and the learning curve decays like $n^{-2/3}$ instead of $1/n$.

\begin{theorem}[The standard form of the log-likelihood ratio]\label{thm:toolkit:empirical}
Assume \cref{ass:toolkit:fundamental}, let $\eps>0$ be small enough, and let $W_\eps^\circ$ be the interior of $\{w\in W:K(w)\le\eps\}$. Then there are a $d$-dimensional real analytic manifold $M$; a proper analytic map $g\colon M\to W_\eps^\circ$ that maps $\{u:K(g(u))\ne0\}$ one-to-one onto $\{w\in W_\eps^\circ:K(w)\ne0\}$; a function $a(x,u)$, analytic in $u$ with values in $L^s(q)$; and an infinitely differentiable function $b>0$, such that in each local coordinate of $M$
\begin{gather*}
K(g(u))=u^{2k},\qquad f(x,g(u))=u^ka(x,u),\\
\varphi(g(u))\,\abs{\det g'(u)}=b(u)\,\abs{u^h},
\end{gather*}
with multi-indices $k$ and $h$ of nonnegative integers, not all $k_j$ zero. Finitely many coordinate neighborhoods, each identified with a cube, cover the preimage of a compact set. Let $\xi_n(u)=n^{-1/2}\sum_{i=1}^n\bigl(u^k-a(X_i,u)\bigr)$.
\begin{enumerate}
\item For every $u$, $\E_X[a(X,u)]=u^k$, so that $\E\,\xi_n(u)=0$, and
\begin{equation}\label{eq:toolkit:standardform}
nK_n(g(u))=nu^{2k}-\sqrt n\,u^k\,\xi_n(u).
\end{equation}
\item As $n\to\infty$, $\xi_n$ converges in law, as a random continuous function on each coordinate cube with the uniform norm, to a centered Gaussian process $\xi$ with covariance $\E[\xi(u)\xi(v)]=\E_X[a(X,u)a(X,v)]-u^kv^k$, and $\E[\sup_u\xi_n(u)^2]\to\E[\sup_u\xi(u)^2]$.
\item If $q$ is realizable by the model, then $\E[\xi_n(u)^2]=\E_X[a(X,u)^2]=2$ at every $u$ with $u^{2k}=0$.
\end{enumerate}
\end{theorem}

The quartet $(M,g,a,b)$ and its formulas are Theorem~1 of \textcite{Watanabe2013}, which states that its proof is given in Theorem~6.1 of the 2009 book \parencite{Watanabe2009}. The argument, as \textcite[Section~4.2]{Watanabe2024} sketches it, combines \cref{thm:toolkit:hironaka} with condition (iv): relatively finite variance is what makes $f(x,g(u))$ divisible by $u^k$. For parts (ii) and (iii), \textcite[p.~881]{Watanabe2013} cites Theorems~6.1--6.3 of the 2009 book, which we know only through that citation; \textcite[Section~4.2, eqs.~(42)--(43)]{Watanabe2024} states the convergence in law, the covariance, and the convergence of $\E\sup_u\xi_n(u)^2$. Restricting to $\{K\le\eps\}$ loses nothing: for $\beta=\beta_0/\log n$ and $0\le r<\frac12$, the integrals of $e^{-n\beta K_n}\varphi$ and of $nK_ne^{-n\beta K_n}\varphi$ over $\{K\ge n^{-r}\}$ are $o_p(e^{-\sqrt n})$ \parencite[Lemma~1]{Watanabe2013}. The book takes (ii) and (iii) on trust; part (i) is a computation.

\begin{proof}[Proof of (i)]
Integrating $f(x,g(u))=u^ka(x,u)$ against $q$ gives $u^{2k}=K(g(u))=u^k\E_X[a(X,u)]$, so $\E_X[a(X,u)]=u^k$ wherever $u^k\ne0$, and by continuity everywhere. Then $nK_n(g(u))=\sum_if(X_i,g(u))=u^k\sum_ia(X_i,u)=u^k(nu^k-\sqrt n\,\xi_n(u))$.
\end{proof}

In the variable $t=nu^{2k}$, identity \eqref{eq:toolkit:standardform} reads $nK_n=t\mp\sqrt t\,\xi_n(u)$, with the sign of $u^k$: the random part of $nK_n$ is of the order of the square root of its mean, uniformly up to the true set. (The number $2$ in (iii) comes from $\E_X[e^{-f}]=1$ in the realizable case, which makes $K=\E_X[e^{-f}-1+f]$ equal to $\frac12\E_X[f^2]$ to second order in $f$.) Substituting \eqref{eq:toolkit:standard} and \eqref{eq:toolkit:standardform} into $\int e^{-nK_n}\varphi$ gives the renormalized posterior, with density in $t$ proportional to $t^{\lambda-1}e^{-t+\sqrt t\,\xi}$ on the coordinates attaining $\lambda$ \parencite[eq.~(48)]{Watanabe2024}, from which the free energy, generalization, WAIC and WBIC theorems of Part III follow.\index{empirical process} \Cref{ch:freeenergy} uses \cref{thm:toolkit:empirical} in its proof sketch of \cref{thm:freeenergy:main}, \cref{ch:generalization,ch:waic,ch:temperature} rely on it for \cref{thm:generalization:main,thm:waic:main,thm:temperature:wbic}, and \cref{ch:llc} uses its local version.

\section{Divergences and modes of convergence}\label{sec:toolkit:probability}

\begin{definition}[Entropy, cross-entropy, and Kullback--Leibler divergence]\label{def:toolkit:kl}
Let $q$ and $p$ be probability densities with respect to a measure on a space $\cX$ (Lebesgue measure for continuous data, counting measure for discrete data), with distributions $Q$ and $P$. The \term{entropy} of $q$ is $S(q)=-\int q\log q$, the \term{cross-entropy} of $p$ relative to $q$ is $-\int q\log p$, and the \term{Kullback--Leibler divergence}\index{Kullback--Leibler divergence} is $\KL(q\,\|\,p)=\int q\log(q/p)\in[0,\infty]$, with $0\log0=0$ and $\KL(q\,\|\,p)=\infty$ when $Q(p=0)>0$. When the entropy is finite, the divergence is the cross-entropy minus the entropy. In the book's notation $L(w)$ is the cross-entropy of $p(\cdot\mid w)$ relative to the truth $q$, $S$ is the entropy of $q$, and in the realizable case $K(w)=L(w)-S=\KL(q\,\|\,p(\cdot\mid w))$. The total variation distance is $\mathrm{TV}(q,p)=\sup_A\abs{Q(A)-P(A)}=\frac12\int\abs{q-p}$.
\begin{enumerate}
\item (Gibbs' inequality.) $\KL(q\,\|\,p)\ge0$, with equality only if $p=q$ almost everywhere.
\item (Pinsker's inequality.) $\mathrm{TV}(q,p)\le\sqrt{\KL(q\,\|\,p)/2}$.
\item $\KL(\Normal(\mu_1,\sigma^2)\,\|\,\Normal(\mu_2,\sigma^2))=(\mu_1-\mu_2)^2/(2\sigma^2)$. Consequently, for a regression $y=f(w,x)+\Normal(0,\sigma^2)$ with inputs $x\sim q(x)$ and the truth $p(\cdot\mid w_0)$, $K(w)=\E_x[(f(w,x)-f(w_0,x))^2]/(2\sigma^2)$.
\end{enumerate}
\end{definition}

\begin{proof}[Proofs]
(i) If the divergence is infinite, then there is nothing to prove. Otherwise $\log x\le x-1$ gives $-\KL(q\,\|\,p)=\int_{q>0}q\log\frac pq\le\int_{q>0}(p-q)=P(q>0)-1\le0$, with equality only if $p=q$ almost everywhere on $\{q>0\}$ and $P(q>0)=1$.
(ii) The plan is to reduce to distributions on two points and then to compare two functions of one variable. Assume the divergence is finite, so that $q=0$ almost everywhere on $\{p=0\}$. For a set $B$ with $P(B)>0$, Jensen's inequality for the convex function $x\log x$ under the density $p\ind_B/P(B)$ gives the log-sum inequality $\int_Bq\log\frac qp\ge Q(B)\log\frac{Q(B)}{P(B)}$. Let $A=\{q>p\}$, $\alpha=Q(A)$ and $\beta=P(A)$, so that $\mathrm{TV}(q,p)=\alpha-\beta\ge0$; if $P(A)$ or $P(A^c)$ is $0$, then so is $\mathrm{TV}(q,p)$. Otherwise the log-sum inequality on $A$ and on $A^c$ gives $\KL(q\,\|\,p)\ge k(\beta)$, where
\[
k(\beta)=\alpha\log\frac\alpha\beta+(1-\alpha)\log\frac{1-\alpha}{1-\beta}.
\]
The function $j(\beta)=k(\beta)-2(\alpha-\beta)^2$ vanishes at $\beta=\alpha$, and for $0<\beta\le\alpha$ its derivative is $(\alpha-\beta)\bigl(4-\frac1{\beta(1-\beta)}\bigr)\le0$, because $\beta(1-\beta)\le\frac14$. So $j\ge0$, and $\KL(q\,\|\,p)\ge2(\alpha-\beta)^2=2\,\mathrm{TV}(q,p)^2$.
(iii) Under $\Normal(\mu_1,\sigma^2)$ the log-ratio $\frac{(x-\mu_2)^2-(x-\mu_1)^2}{2\sigma^2}$ has mean $\frac{(\mu_1-\mu_2)^2}{2\sigma^2}$. For the regression, the factor $q(x)$ cancels in the ratio of joint densities, and the divergence is the mean over $x$ of the divergence between the two conditional normal densities.
\end{proof}

When the truth is realizable and the model is smooth, the Hessian of $K$ at a true parameter is the Fisher information \parencite[p.~872]{Watanabe2013}. \Cref{ch:regular,ch:singular} identify $K$ with a divergence and with half the mean squared difference of the regression functions, and \cref{ch:generalization,ch:waic} measure generalization by the divergence from the truth to the predictive distribution.\sidenote{Pinsker's inequality is named after \textcite{Pinsker1964}; the form with the constant $\frac12$ is in \textcite{Kullback1967}.}\index{Pinsker's inequality}

\begin{definition}[Modes of convergence]\label{def:toolkit:convergence}
Let $Y_n$ and $Y$ be random vectors in $\R^k$. The sequence $Y_n$ converges \term{in probability} to $Y$ (on a common probability space) if $\Prob(\abs{Y_n-Y}>\eta)\to0$ for every $\eta>0$, and \term{in law} if $\E h(Y_n)\to\E h(Y)$ for every bounded continuous $h$; for real random variables, convergence in law is convergence of the distribution functions at each continuity point of the limit. Random continuous functions on a compact set converge in law if the same holds for every bounded $h$ that is continuous for the uniform norm. The sequence is $O_p(1)$ if for every $\eta>0$ some $C$ has $\Prob(\abs{Y_n}>C)\le\eta$ for all $n$; $Y_n=O_p(a_n)$ means $Y_n/a_n=O_p(1)$, and $Y_n=o_p(a_n)$ means $Y_n/a_n\to0$ in probability.
\begin{enumerate}
\item Convergence in probability implies convergence in law; convergence in law to a constant implies convergence in probability; and a sequence that converges in law is $O_p(1)$.
\item (Continuous mapping.) If $Y_n\to Y$ in law and $\Phi$ is continuous, then $\Phi(Y_n)\to\Phi(Y)$ in law.
\item (Slutsky's lemma.) If $Y_n\to Y$ in law and $Z_n\to c$ in probability for a constant $c$, then $(Y_n,Z_n)\to(Y,c)$ in law. Hence $Y_n+Z_n\to Y+c$ and $Z_nY_n\to cY$ in law, and, for real $Z_n$ and $c\ne0$, $Y_n/Z_n\to Y/c$.
\item $O_p(1)\,o_p(1)=o_p(1)$ and $O_p(1)+o_p(1)=O_p(1)$; and if $Y_n\to Y$ in law and $R_n=o_p(1)$, then $Y_n+R_n\to Y$ in law.
\end{enumerate}
\end{definition}

The book uses the \term{portmanteau lemma}: convergence in law holds as soon as $\E h(Y_n)\to\E h(Y)$ for every bounded Lipschitz $h$ \parencite[Lemma~2.2]{vanderVaart1998}. With it, the first claim of (i) follows from $\abs{\E h(Y_n)-\E h(Y)}\le L\eta+2B\,\Prob(\abs{Y_n-Y}>\eta)$ for $h$ bounded by $B$ with Lipschitz constant $L$, and the second from $\Prob(\abs{Y_n-c}>\eta)\le\E\min(1,\abs{Y_n-c}/\eta)$; the third claim and (iv) are proved in \textcite[Chapter~2]{vanderVaart1998}; and (ii) holds because $h\circ\Phi$ is bounded and continuous whenever $h$ is.

\begin{proof}[Proof of Slutsky's lemma]
The plan is to check the portmanteau criterion for the pairs and then to apply continuous maps. Let $h$ be bounded by $B$ and Lipschitz with constant $L$. For every $\eta>0$,
\[
\abs{\E h(Y_n,Z_n)-\E h(Y_n,c)}\le\E\min\bigl(L\abs{Z_n-c},2B\bigr)\le L\eta+2B\,\Prob(\abs{Z_n-c}>\eta),
\]
whose upper limit is at most $L\eta$; and $\E h(Y_n,c)\to\E h(Y,c)$, because $y\mapsto h(y,c)$ is bounded and continuous. Letting $\eta\to0$ gives $\E h(Y_n,Z_n)\to\E h(Y,c)$, so $(Y_n,Z_n)\to(Y,c)$ in law, and the sum and product follow from (ii). For the quotient with $c>0$, the function $\Phi(y,z)=y/\max(z,c/2)$ is continuous, so $\Phi(Y_n,Z_n)\to Y/c$ in law; and $Y_n/Z_n=\Phi(Y_n,Z_n)$ on the event $\{Z_n>c/2\}$, whose probability tends to $1$, so the difference tends to $0$ in probability and the sum rule gives $Y_n/Z_n\to Y/c$. The case $c<0$ is the same with $\min(z,c/2)$.
\end{proof}

The book's asymptotic statements about random quantities hold in law, in probability, or in expectation over the data; the proofs in \cref{ch:regular,ch:coefficient,ch:generalization,ch:phases} use Slutsky's lemma to discard $o_p(1)$ terms.\index{Slutsky's lemma}\index{convergence in law}

\section{Samplers and estimators}\label{sec:toolkit:mcmc}

The book's measurements are made on models with few parameters, so most of them use deterministic quadrature; Markov chain Monte Carlo appears where quadrature is impossible, and as a check of the samplers on models where it is possible. This section describes what the kit \texttt{figures/\allowbreak scripts/\allowbreak sltkit.py} implements and what can go wrong. The energy is $E(w)=nL_n(w)$, the prior is uniform on a box unless stated, and the tempered posterior at inverse temperature $\beta$ has density proportional to $\varphi e^{-\beta E}$.

\paragraph{Quadrature on a grid.} For $d\le2$, \texttt{free\_\allowbreak energy\_\allowbreak grid} evaluates $E$ at the midpoints $w_j$ of a regular grid on the box, with cell volume $\abs\Delta$, and computes
\[
F_n(\beta)\approx-\log\sum_j\frac{\abs\Delta}{\vol(W)}e^{-\beta E(w_j)}
\]
as a log-sum of exponentials from which the largest exponent is subtracted first; $e^{-E}$ itself underflows once $n$ is in the hundreds. The function \texttt{tempered\_\allowbreak mean\_\allowbreak grid} averages with the weights $p_j\propto e^{-\beta E(w_j)}$, and \texttt{sample\_\allowbreak grid} draws a cell with probability $p_j$ and a uniform point inside it, an exact sample from the piecewise-constant approximation of the posterior. The midpoint rule is accurate when the integrand varies smoothly on the scale of the cells, and it degrades once the posterior is much narrower than a cell. For the product model the posterior hugs the cross in a band of width about $n^{-1/2}/\abs a$; on one data set at each $n$, $1600$ points per side of $[-1,1]$ reproduce the exact one-dimensional free energy below to $10^{-7}$ up to $n=10^6$, but are off by $2\times10^{-3}$ at $n=10^7$ and by $0.25$ at $n=10^8$, where the band is a fraction of a cell wide. A script should check its grid by doubling it.

\paragraph{One coefficient.} When the energy depends on $w$ only through a scalar $c=g(w)$ and the noise is Gaussian, it is quadratic in $c$: $E=\hat E+\frac{S_{xx}}{2\sigma^2}(c-\hat c)^2$, with $S_{xx}=\sum_ix_i^2$, $\hat c=\sum_ix_iy_i/S_{xx}$, and $\hat E$ the minimum over all real $c$. Then $F_n(\beta)=\beta\hat E-\log\int e^{-\beta S_{xx}(c-\hat c)^2/(2\sigma^2)}\rho(c)\dd c$, with $\rho$ the density of $g(w)$ under the prior: $(\log\frac1{\abs c})^{L-1}/(2(L-1)!)$ for $c=a_1\cdots a_L$ with independent uniform factors on $[-1,1]$, and $(\abs c^{-1/2}-1)/2$ for $c=ab^2$. The functions \texttt{coefficient\_\allowbreak free\_\allowbreak energy} and \texttt{free\_\allowbreak energy\_\allowbreak chain} integrate over $\hat c\pm40$ posterior standard deviations, split at the singular points of $\rho$, by adaptive quadrature to relative accuracy $10^{-10}$.

\paragraph{Exact draws for $y=abx$.} Under the uniform prior on $[-1,1]^2$, the change of variables $(a,b)\mapsto(a,c)$ with $c=ab$ has Jacobian $1/\abs a$, so given $c$ the density of $a$ is proportional to $1/\abs a$ on $\abs c\le\abs a\le1$: $\log\abs a$ is uniform on $[\log\abs c,0]$. The function \texttt{sample\_\allowbreak product\_\allowbreak posterior} draws $c$ from its one-dimensional posterior by inverting a distribution function tabulated on about $8\times10^4$ points, dense near $0$, then $\log\abs a$, a random sign, and $b=c/a$; the draws are exact up to the interpolation of that table.

\paragraph{Random-walk Metropolis.} The function \texttt{metropolis} runs many chains at once.\sidenote{The algorithm is that of \textcite{MetropolisEtAl1953}.} From the current point $w$ it proposes $w'=w+h\,\xi$ with $\xi\sim\Normal(0,I_d)$ and accepts with probability $\min\{1,\pi(w')/\pi(w)\}$, where $\pi\propto\varphi e^{-\beta E}$; the scripts' log-densities are $-\infty$ outside the box, so a proposal outside it is rejected. Since $\pi(w)\min\{1,\pi(w')/\pi(w)\}=\min\{\pi(w),\pi(w')\}$ is symmetric in $(w,w')$, the chain satisfies detailed balance and leaves $\pi$ invariant. During the first steps, which are discarded, each chain multiplies its $h$ by $e^{0.05(1-0.3)}$ after an acceptance and by $e^{-0.05\cdot0.3}$ after a rejection (for the default target $0.3$), a recursion whose mean drift vanishes at acceptance rate $0.3$; then $h$ is frozen, so the kept draws come from an ordinary Metropolis chain. The rate $0.3$ lies above $0.234$, the optimal acceptance rate for densities with many independent coordinates \parencite{RobertsEtAl1997}, because the book's posteriors have few. For a positive continuous density on a box the chain is irreducible and aperiodic, and its averages converge to posterior means \parencite{Tierney1994}, but slowly on a singular posterior: across an arm of the cross its width is about $n^{-1/2}/\abs a$, along the arm it reaches the edge of the box (\cref{thm:singular:scales}), and a step size accepted at the narrow end diffuses only slowly along $\log\abs a$. Separated regions, such as the phases of \cref{ch:phases}, may not communicate at all. The defence is to run chains from dispersed starting points and compare them, for instance with the rank-normalized $\hat R$ of \textcite{VehtariEtAl2021}, and to validate every Metropolis measurement on a running model whose answer is known.\index{Metropolis algorithm}

\paragraph{Hamiltonian Monte Carlo and tempering.} The kit implements neither. Hamiltonian Monte Carlo \parencite{Neal2011} adds a Gaussian momentum, follows Hamilton's equations for the energy $-\log\pi$ with the leapfrog scheme, and accepts by a Metropolis test; it scales better with dimension, but its step size is also set by the narrowest direction of the posterior. Tempering runs chains at a ladder of inverse temperatures, which helps when regions that do not communicate at $\beta=1$ do at smaller $\beta$.

\paragraph{Thermodynamic integration.} With a normalized prior, the tempered free energy $F_n(\beta)=-\log\int\varphi e^{-\beta E}$ vanishes at $\beta=0$, and differentiation under the integral sign gives
\[
F_n'(\beta)=\Ewb{nL_n(w)},\qquad F_n''(\beta)=-\Var^\beta_w\bigl(nL_n(w)\bigr)\le0,
\]
so that $F_n=\int_0^1\Ewb{nL_n(w)}\dd\beta$. The function \texttt{thermodynamic\_\allowbreak integration} applies the trapezoid rule on the ladder $\beta_k=(k/K)^5$, $k=0,\dots,K$, of \texttt{power\_\allowbreak ladder}, which is dense near $\beta=0$, where the integrand changes fastest: its excess over $nL_n(w_0)$ is of order $n$ at $\beta=0$, where the tempered posterior is the prior, and about $\lambda/\beta$ once $n\beta$ is large (\cref{ch:temperature}), so the range $1/n\ll\beta\le1$ contributes the term $\lambda\log n$. This is the power posterior method of \textcite{FrielPettitt2008}. Two errors add up: the discretization of the integral, largest near $\beta=0$, and the sampling error of each rung, largest where the tempered posterior is wide or changes phase.\index{thermodynamic integration}

\paragraph{WAIC and WBIC.} From draws $w_1,\dots,w_S$ of the posterior at $\beta=1$, \texttt{waic} computes
\begin{gather*}
T_n=-\frac1n\sum_{i=1}^n\log\Bigl(\frac1S\sum_{s=1}^Sp(X_i\mid w_s)\Bigr),\qquad V_n=\sum_{i=1}^n\widehat\Var_s\bigl[\log p(X_i\mid w_s)\bigr],\\
\WAIC=T_n+V_n/n,
\end{gather*}
with the inner mean taken as a log-sum of exponentials and $\widehat\Var_s$ the sample variance over the draws, with divisor $S-1$. The formula assumes exact draws at $\beta=1$, and correlated or few draws add Monte Carlo error. The book's $\WAIC$ is a loss per sample; the deviance scale of \textcite{GelmanHwangVehtari2014}, used by common software, is $2n$ times it. The function \texttt{wbic} averages the energy over draws at $\beta=1/\log n$ \parencite{Watanabe2013}; \cref{thm:temperature:wbic} says how close the result is to $F_n$.

\paragraph{Two temperatures.} Since $\Ewb{nL_n}$ is about $nL_n(w_0)+\lambda/\beta$, the difference of two tempered means cancels the unknown $nL_n(w_0)$:
\[
\hat\lambda=\frac{\E^{\beta_1}_w[nL_n(w)]-\E^{\beta_2}_w[nL_n(w)]}{1/\beta_1-1/\beta_2},\qquad\beta_i=\frac{\beta_{0i}}{\log n},
\]
which converges to $\lambda$ in probability \parencite[Corollary~3]{Watanabe2013} and is computed by \texttt{lambda\_\allowbreak two\_\allowbreak temperatures}; the second mean can be obtained from draws at the first temperature by reweighting with $e^{-(\beta_2-\beta_1)nL_n}$ \parencite[eq.~(20)]{Watanabe2013}. For $m>1$ the next term of the tempered mean, $-(m-1)/(\beta\log(n\beta))$ in the population version of \cref{prop:temperature:mean}, biases the estimate by about $-(m-1)/\log n$; it is the $\beta$-derivative of the term $-(m-1)\log\log(n\beta)$ of the free energy, and \textcite{Watanabe2013} notes that this term affected his estimates for reduced-rank regression with $m=2$.

\paragraph{Volumes.} The function \texttt{volume\_\allowbreak fraction} estimates $V(\eps)$ by the fraction of $N$ uniform draws with $K<\eps$, with Clopper--Pearson intervals. The relative error is about $(NV(\eps))^{-1/2}$, so $N$ must be many times $1/V(\eps)\approx\eps^{-\lambda}$. The least-squares slope of $\log V$ against $\log\eps$, from \texttt{loglog\_\allowbreak slope}, estimates $\lambda$ only when $m=1$: for $V\sim C\eps^\lambda(\log\frac1\eps)^{m-1}$ the local slope is $\lambda-(m-1)/\log\frac1\eps$.

Most chapters' scripts use the deterministic methods of this section. Metropolis with the two-temperature estimator, the second mean obtained by reweighting, serves \cref{ch:catalogue,ch:temperature}; thermodynamic integration and $\WBIC$ serve \cref{ch:temperature}, whose scripts compute them by quadrature in their own code; and \texttt{waic} serves \cref{ch:waic}.

\section{Stochastic-gradient Langevin dynamics}\label{sec:toolkit:sgld}

The posterior of a trained network cannot be computed on a grid, and Metropolis chains on millions of parameters are too slow. Part IV therefore estimates the local learning coefficient from a localized tempered posterior sampled by stochastic-gradient Langevin dynamics, following \textcite{LauEtAl2025}. For a parameter $w^*$, an inverse temperature $\beta$, and a localization strength $\gamma>0$, the \term{localized tempered posterior}\index{localized posterior} is the probability density on $\R^d$
\[
p(w\mid w^*,\beta,\gamma)\propto\exp\bigl(-n\beta L_n(w)-\tfrac\gamma2\abs{w-w^*}^2\bigr),
\]
in which a Gaussian factor centered at $w^*$ replaces the prior, and the estimator of the local learning coefficient is
\[
\hat\lambda(w^*)=n\beta^*\bigl(\E_{w\mid w^*,\beta^*,\gamma}[L_n(w)]-L_n(w^*)\bigr),\qquad\beta^*=1/\log n,
\]
the WBIC idea applied near one point \parencite[eqs.~(10) and (12)]{LauEtAl2025}. Consider its population version, at a true parameter $w^*$ and with $K$ in place of $K_n$. If the volume of $\{K<\eps\}$ under the Gaussian weight behaves like $C\eps^{\tilde\lambda}(\log\frac1\eps)^{\tilde m-1}$, then part (i) of \cref{thm:toolkit:karamata}, with that weight in place of the prior, shows that the estimate tends to $\tilde\lambda$ as $n\beta\to\infty$ with $\gamma$ fixed. Because the weight is positive on all of $\R^d$, $\tilde\lambda$ is set by the most singular true parameters anywhere, not only near $w^*$; \textcite{LauEtAl2025} conjecture that their localized coefficient is the local learning coefficient under conditions on $\gamma$, and \cref{ch:llc} discusses what is known.

The Langevin diffusion $\dd w=-\frac12\nabla U(w)\dd t+\dd B_t$, with $B$ a standard Brownian motion, has stationary density proportional to $e^{-U}$, and its Euler discretization with step $\epsilon$ is $w_{t+1}=w_t-\frac\epsilon2\nabla U(w_t)+\eta_t$ with independent $\eta_t\sim\Normal(0,\epsilon I_d)$. \term{Stochastic-gradient Langevin dynamics}\index{stochastic-gradient Langevin dynamics} replaces $\nabla U$ by an unbiased estimate computed on a minibatch.\sidenote{\Textcite{WellingTeh2011} introduced the method, with step sizes that decrease to zero.} For $U(w)=n\beta L_n(w)+\frac\gamma2\abs{w-w^*}^2$ and a minibatch $B_t$ of $b$ indices drawn at random at each step, the update is
\begin{equation}\label{eq:toolkit:sgld}
w_{t+1}=w_t+\frac\epsilon2\Bigl(\frac{n\beta}b\sum_{i\in B_t}\nabla\log p(X_i\mid w_t)+\gamma(w^*-w_t)\Bigr)+\eta_t,
\end{equation}
which is the update of \textcite[Section~4.4]{LauEtAl2025} in the notation of this book (their minibatch size $m$ is our $b$). The estimate replaces the posterior mean by the average of $L_n(w_t)$ over the steps kept after a burn-in; the authors recycle the minibatch losses computed for the gradients as unbiased estimates of $L_n(w_t)$.

The method has five known weaknesses. First, with a constant step the chain does not sample the target: decreasing steps give consistency and a central limit theorem \parencite{TehEtAl2016}, but practice uses constant steps, and with the usual step proportional to $1/n$ the invariant law of SGLD departs from the posterior as the data grow and comes to resemble that of stochastic gradient descent, as \textcite{BrosseEtAl2018} prove for strongly log-concave posteriors. Second, too small a step does not equilibrate and underestimates, and too large a step is unstable and overestimates; \textcite[Appendix~H]{LauEtAl2025} tune $\epsilon$ with a Metropolis-adjusted acceptance diagnostic aimed at $0.9$--$0.95$. Third, too large a $\gamma$ overwhelms the likelihood, and the Gaussian localizer, which is not compactly supported, lies outside the setting of the theory, as \textcite{HitchcockHoogland2025} note in a preprint. Fourth, if $w^*$ is not near a local minimum, then the chain finds lower loss and the estimate can be negative. Fifth, the global convergence guarantees for stochastic-gradient samplers assume conditions that degenerate landscapes do not satisfy \parencite{HitchcockHoogland2025}: as of 2026 we know of no theorem that covers the chains used to estimate the local learning coefficient, and their support is empirical, from known values for deep linear networks and a comparison with the Metropolis-adjusted Langevin algorithm on a small network \parencite{LauEtAl2025}, and from the same preprint's benchmark of local samplers on deep linear networks with known local values. Conventions for $n\beta$ and $\gamma$ differ between papers, so estimates are comparable only within one convention.

\Cref{ch:llc} defines the local learning coefficient and estimates it with \eqref{eq:toolkit:sgld}, \cref{ch:stages} follows the estimate along training, and \cref{ch:sgd} compares Langevin dynamics with the stochastic gradient descent it modifies.

\section{Intervals}\label{sec:toolkit:intervals}

Most error bars in the book are $95\%$ intervals of one of the three kinds below; a few are normal-approximation intervals, an estimate plus or minus $1.96$ standard errors, or intervals for a fitted slope or a correlation. Every caption says which, and over how many independent replicates.

\paragraph{Student-$t$ intervals.} For $R$ independent replicates $v_1,\dots,v_R$ of a quantity, such as a free energy on $R$ data sets, with mean $\bar v$ and sample variance $s^2=\sum_r(v_r-\bar v)^2/(R-1)$, \texttt{mean\_ci} reports $\bar v\pm t_{R-1}s/\sqrt R$, with $t_{R-1}$ the $0.975$ quantile of Student's $t$ distribution with $R-1$ degrees of freedom. If the replicates are normal, then the coverage is exactly $95\%$, because $\sqrt R(\bar v-\E v)/s$ has that distribution; if they only have a finite variance, then the coverage tends to $95\%$ as $R$ grows, by the central limit theorem and Slutsky's lemma (\cref{def:toolkit:convergence}). For skewed or heavy-tailed replicates and small $R$ the coverage can miss $95\%$, and the interval is wrong for the dependent draws of one Markov chain, which are not replicates.

\paragraph{Clopper--Pearson intervals.}\index{Clopper--Pearson interval} For $J$ successes in $N$ independent trials with success probability $\theta$, such as $J$ uniform draws in $\{K(w)<\eps\}$, \texttt{clopper\_\allowbreak pearson} reports $[\theta_-(J),\theta_+(J)]$ \parencite{ClopperPearson1934}. Here $\theta_-(k)$ is the $0.025$ quantile of the Beta distribution with parameters $(k,N-k+1)$, or $0$ if $k=0$, and $\theta_+(k)$ is the $0.975$ quantile of the Beta distribution with parameters $(k+1,N-k)$, or $1$ if $k=N$.

\begin{proposition}\label{prop:toolkit:clopper}
For every $\theta\in[0,1]$, $\Prob_\theta\bigl(\theta_-(J)\le\theta\le\theta_+(J)\bigr)\ge0.95$.
\end{proposition}

\begin{proof}
The plan is to express binomial tails as Beta probabilities and then to use monotonicity. For $1\le k\le N$, $\Prob_\theta(J\ge k)=\Prob(B_k\le\theta)$, where $B_k$ has the Beta distribution with parameters $(k,N-k+1)$. Indeed, if $U_1,\dots,U_N$ are independent and uniform on $[0,1]$, then the number of $U_i\le\theta$ has the law of $J$, and it is at least $k$ exactly when the $k$-th smallest $U_i$ is at most $\theta$; that order statistic has density $\frac{N!}{(k-1)!(N-k)!}u^{k-1}(1-u)^{N-k}$, which is the Beta density with these parameters. The quantiles $\theta_+(k)$ increase with $k$. If some $k$ has $\theta_+(k)<\theta$, then let $k^*$ be the largest; it satisfies $k^*<N$, the event $\{\theta_+(J)<\theta\}$ is $\{J\le k^*\}$, and
\[
\Prob_\theta(J\le k^*)=\Prob(B_{k^*+1}>\theta)\le\Prob\bigl(B_{k^*+1}>\theta_+(k^*)\bigr)=0.025,
\]
because $\theta_+(k^*)$ is the $0.975$ quantile of $B_{k^*+1}$ and $\theta>\theta_+(k^*)$. In the same way the event $\{\theta_-(J)>\theta\}$, if it is not empty, is $\{J\ge k_*\}$ for the smallest $k_*\ge1$ with $\theta_-(k_*)>\theta$, and $\Prob_\theta(J\ge k_*)=\Prob(B_{k_*}\le\theta)\le\Prob(B_{k_*}\le\theta_-(k_*))=0.025$. The two failures together have probability at most $0.05$.
\end{proof}

The interval is conservative: the infimum of its coverage over $\theta$, approached just outside an end point of one of the intervals, is $96.10\%$ for $N=10$ and $95.002\%$ for $N=200$. \Cref{ch:volume,ch:newton} use it for Monte Carlo volume fractions, and other chapters for the frequencies of events.

\paragraph{Bootstrap intervals.}\index{bootstrap} For a statistic that is not a mean, such as an exponent fitted to Monte Carlo volume fractions at several levels $\eps$ in \cref{ch:volume}, a script that reports a bootstrap interval resamples the $R$ independent replicates (there, the uniform draws, whose resampling amounts to multinomial counts) with replacement $B$ times, recomputes the statistic on each resample, and reports the $2.5\%$ and $97.5\%$ quantiles of the $B$ values: the percentile interval of \textcite{Efron1979}. Its coverage tends to $95\%$ for smooth functions of means of independent replicates. It can fail for nonsmooth statistics and for parameters on the boundary of the parameter space, where the bootstrap is inconsistent \parencite{Andrews2000}; estimators at a singular point, whose limit laws change with the truth, are of the same nonregular kind, so a bootstrap interval for such an estimator needs a check by simulation. Resampling the correlated draws of one chain understates the error.

\notesandrefs

\textcite{Wong2001} treats Laplace's method and the Mellin transform technique for integrals. \textcite{BinghamEtAl1987} is the standard reference on regular variation; the proof of part (i) of \cref{thm:toolkit:karamata} was written for this book, as the usual dominated-convergence argument made elementary by the explicit form of $V$. \textcite{KrantzParks2002} develop real analytic functions, and \textcite{vanderVaart1998} the modes of convergence. Hironaka's theorem, Watanabe's standard representation and the theorem on the empirical process are taken from the sources named where they are stated, and where we know a source only through another's citation we say so; WAIC is from \textcite{Watanabe2010}.

\paragraph{Translating the literature.} Watanabe's 2009 book writes $f(x,w)=\log\{q(x)/p(x\mid w)\}$ in the realizable case, and his later papers $f(x,w)=\log\{p_0(x)/p(x\mid w)\}$; their $K$ and $K_n$ are ours, and \textcite{Watanabe2024} writes $\theta$ and $\pi(\theta)$ for our $w$ and $\varphi(w)$. \textcite{BinghamEtAl1987} write $U$ for our $V$ and $\hat U(s)$ for our $Z(n)$, and state Karamata's theorem at infinity rather than at $0$. \textcite{LauEtAl2025} write $m$ for the minibatch size that we call $b$ (our $m$ is the multiplicity) and $\beta^*$ for $1/\log n$. \textcite{GelmanHwangVehtari2014} report WAIC on the deviance scale, $2n$ times ours.

\printchapterbib
\end{refsection}
